%! Author = Joey Dufourd
%! Date = 29/10/2025

\documentclass[11pt,a4paper]{article}
\usepackage[margin=1.8cm]{geometry}
\usepackage[T1]{fontenc}
\usepackage[utf8]{inputenc}
\usepackage{lmodern}
\usepackage{setspace}
\usepackage{enumitem}
\usepackage{titlesec}
\usepackage[hidelinks]{hyperref}
\usepackage{array}
\usepackage{multicol}
\usepackage{tabularx}
\usepackage{parskip}
\usepackage{fontawesome5}
\usepackage{graphicx}

% ---- Compact section titles
\titleformat{\section}{\LARGE\bfseries}{}{0pt}{}
\titlespacing*{\section}{0pt}{10pt plus 1pt minus 1pt}{6pt}


\setlist[itemize]{leftmargin=1.2em,itemsep=0.5pt,topsep=1pt}
\newcommand{\sep}{\;\textbar\;}
\let\oldlabelitemi\labelitemi
% \renewcommand{\labelitemi}{\faCode} \renewcommand{\labelitemi}{\faMicrochip} \renewcommand{\labelitemi}{\faCogs} \renewcommand{\labelitemi}{\faTerminal}
% pdflatex -output-directory=assets CV-JoeyDufourd.tex

% ---- Metadata
\hypersetup{
  pdftitle={Curriculum Vitae - Joey Dufourd},
  pdfauthor={Joey Dufourd <joey.dufourd@gmail.com>},
  pdfsubject={Postdoctoral Application},
  pdfkeywords={Computational Biology, Bioinformatics, Chromatin, Telomeres, Multi-omics, Machine Learning, Image Analysis, Python, R},
  pdfproducer={LaTeX with hyperref},
  pdfcreator={pdflatex}
}



\begin{document}

% =================== HEADER ===================
\begin{center}
{\LARGE \textbf{Joey Dufourd, Ph.D.}}\\[4pt]
\vspace{2pt}
\rule{0.4\textwidth}{0.5pt}\\[8pt]
\vspace{2pt}
\faEnvelope~\href{mailto:joey.dufourd@gmail.com}{joey.dufourd@gmail.com} \sep \faPhone~+33 6 41 28 00 44 \sep
\faGithub~\href{https://github.com/joeydufourd}{github.com/joeydufourd} \sep \faGlobe~\href{https://joeydufourd.github.io/}{https://joeydufourd.github.io} \sep
\faOrcid~\href{https://orcid.org/0000-0002-4138-3672}{0000-0002-4138-3672} \sep
\faLinkedin~\href{https://linkedin.com/in/joey-dufourd}{linkedin.com/in/joey-dufourd}
\\[6pt]
\vspace{2pt}
\rule{0.8\textwidth}{0.5pt}\\[8pt]
\vspace{2pt}
\textbf{Research interests:} \textit{Life Science \sep Computational Biology \& BioInformatics \sep Chromatin \& Telomeres \sep Multi-omics \sep Machine Learning \sep Computer Vision \sep Big data \& Data Science}
\end{center}

% =================== SUMMARY ===================
\vspace{1.5pt}
\noindent\rule{\textwidth}{0.4pt}
\section*{Profile}
Molecular biologist with computational expertise, passionate about understanding life through science.
My PhD on telomere biology combined advanced molecular and imaging approaches (3D-SIM, proteomics) with custom computational pipelines in Python and R for image analysis, statistical modeling, computer vision and data integration.
Proficient in Bash, Python, and R, I leverage computational approaches to develop reproducible analytical workflows and explore complex biological systems.
I am convinced that integrating the strengths of wet-lab experimentation with the power of computational approaches will accelerate discoveries and enable transformative advances in biomedical research.\\
\noindent\rule{\textwidth}{0.4pt}

% =================== RESEARCH EXPERIENCE ===================
\section*{Research Experience}
\textbf{PhD Student --- Repeated sequences chromatin \& Computer vision} \hfill 10/2021--12/2025\hspace{0.3em}\\
Institut de Génétique Humaine (CNRS), Montpellier, France \hfill \textit{\textbf{Supervisors:} J.\ Déjardin, M.\ Tardat}
\begin{itemize}
  \item \textbf{Discovered and characterized TELS1}, a telomere binding \& mouse embryonic stem cell specific protein that stabilizes \textit{t-loops} independently of TRF2.
  \item \textbf{Combined} PICh/IP-MS, ChIP, FISH (nuclei \& metaphases), microinjection of mouse embryos, CRISPR-Cas9, 3D-SIM super-resolution microscopy and other approches to investigate the role of TELS1 at telomeres.
  \item Built \textbf{automation \& AI pipelines} (computer vision) to automate \textit{t-loop} quantification and microscopy analysis (CellProfiler).
  \item Integrated imaging, proteomics, statistics and molecular data using \textbf{Python \& R}.
  \item \textbf{Conducted statistical modeling and experimental design} (ANOVA and multiple-testing aware post-hoc).
  \item \textbf{Co-founded and coordinate} the IGH Young Scientist Club (YSC), a monthly forum promoting collaboration and knowledge sharing among early-career researchers.
  \item Delivered \textbf{seminars and workshops} on computational biology, statistics, and artificial intelligence for students and researchers at IGH.
  \item \textbf{Taught} undergraduate courses in genetics and epigenetics.
\renewcommand{\labelitemi}{\faTerminal}
\renewcommand{\labelitemi}{\raisebox{0.2ex}{\scalebox{0.7}{\faTerminal}}}
  \item \textbf{Tech stack:} Python (TensorFlow/PyTorch, scikit-learn, pandas, NumPy, Matplotlib,\ldots), R (tidyverse, ggplot2, Bioconductor, \ldots), Bash, 3D-SIM image processing, AlphaFold modeling, Git, statistical modeling (ANOVA, clustering, PCA), CellProfiler, ImageJ/Fiji.
\end{itemize}
\renewcommand{\labelitemi}{\oldlabelitemi}


\textbf{M2 Internship --- Repeated Sequences Chromatin} \hfill 02/2021--08/2021\hspace{0.3em}\\
Institut de Génétique Humaine (CNRS), Montpellier, France
\begin{itemize}
  \item Investigated telomeric chromatin structure and \textit{t-loop} stability in \textbf{mouse embryonic stem cell and embryo-like structures}.
  \item \textbf{Analyzed sequencing data}, performed R-based statistical analyses, and quantified images using CellProfiler (normality tests, ANOVA, Bonferroni-corrected tests).
  \item \textbf{Integrated experimental and computational approaches} to identify novel regulators of telomere structure and chromatin dynamics.
\end{itemize}

\textbf{M1 Internship --- Stanford University (Van Rechem Lab)} \hfill 02/2020--08/2020\hspace{0.3em}\\
Stanford Medicine, USA
\begin{itemize}
  \item Characterized non-canonical cytoplasmic functions of \textbf{PRMT5} in cancer biology; applied RNAi, polysome profiling (sucrose gradients), and confocal imaging.
  \item Managed and performed experiments across $>$10 different cancer cell lines in parallel, \textbf{ensuring reproducibility and consistency of results}.
\end{itemize}

\textbf{iGEM Team Leader \& Project Inventor --- Montpellier} \hfill 11/2018--11/2019\hspace{0.3em}\\
\emph{International Genetically Engineered Machine (iGEM) Competition --- Gold Medal}\\
\emph{Project:} \href{https://2019.igem.org/Team:Montpellier}{\textbf{KARMA}} --- an innovative molecular system for targeted protein degradation.
\begin{itemize}
  \item Conceived and led the \textbf{KARMA project}, a modular tool combining specific recognition antibodies with a targeted degradation protease to selectively remove proteins, with potential applications in cancer and antibiotic resistance.
  \item \textbf{Supervised and coordinated} a multidisciplinary team of 14 students (biology, bioinformatics, and engineering) through project design, cloning, modeling, and validation phases.
  \item Built and maintained the \textbf{project website and wiki}, ensuring clear and accessible communication of objectives and results (\href{https://2019.igem.org/Team:Montpellier}{igem.org/Team:Montpellier}).
  \item \textbf{Led public engagement initiatives} (\textit{Carascience}) to present the project and basic biology to schools and the general public.
  \item Strengthened skills in \textbf{project management, communication, synthetic biology, bioinformatic coding, statistical analysis, and web development} within an international competition framework.
\end{itemize}

\textbf{Research Intern --- CNRS (DIMNP \& IGF, Montpellier)} \hfill 10/2018--12/2018\hspace{0.3em}\\
\emph{Supervised by Prof. Ovidiu Radulescu (DIMNP)} --- Project on mathematical modeling of HIV latency and apoptosis mechanisms.
\begin{itemize}
  \item Designed and analyzed \textbf{multiscale computational models} of HIV latency to explore viral persistence and reactivation dynamics across molecular, cellular, and population scales.
  \item Applied \textbf{stochastic simulations (Gillespie algorithm)} and differential modeling to investigate latency transitions; developed reusable scripts later adapted to model telomere state dynamics (\href{https://github.com/joeydufourd/telomere-dynamics-gillespie}{\texttt{github.com/joeydufourd/telomere-dynamics-gillespie}}).
  \item Strengthened skills in \textbf{quantitative biology, dynamical systems, and data analysis} applied to virology and cell fate modeling.
\end{itemize}

\vspace{2pt}
\noindent\rule{\textwidth}{0.4pt}
% =================== PUBLICATIONS ===================
\section*{Publications \& Manuscripts}
\renewcommand{\labelitemi}{\faBook}
\begin{itemize}
  \item \textbf{Dufourd J.} et al. \textit{TELS1 stabilizes telomeric \textit{t}-loops independently of TRF2 in pluripotent cells.} (2025).
  \item \textbf{Dufourd J.} \textit{AI-assisted image analysis for telomeric structure detection and quantification.} In preparation (2025?).
\end{itemize}
\renewcommand{\labelitemi}{\oldlabelitemi}
\vspace{2pt}
\noindent\rule{\textwidth}{0.4pt}
% =================== Projects Highlights ===================
\section*{Projects Highlights}
\begin{itemize}
  \item \textbf{AI-assisted image analysis for telomere structure} — Developed a YOLO-based computer vision pipeline for \textit{t-loop} \& line detection (2025).
  \item \textbf{Stochastic simulations (Gillespie algorithm)} --- Stochastic simulation of telomere state transitions using the Gillespie SSA; explores dynamics of \textit{t-loop} formation/processing and end-protection; includes a reproducible script for parameter sweeps and seeded runs. \emph{(\faGithub~\href{https://github.com/joeydufourd/telomere-dynamics-gillespie}{joeydufourd/telomere-dynamics-gillespie})}
  \item \textbf{FIJI-Macros (ImageJ)} --- Open-source macro suite to automate multichannel microscopy workflows (batch conversion, channel alignment, montage generation,\ldots). \emph{(\faGithub~\href{https://github.com/joeydufourd/FIJI-Macros}{joeydufourd/FIJI-Macros})}
  \item \textbf{Discovery of TELS1} — Identified a pluripotency-specific \textit{t-loop} stabilizing protein through unbiased \& high-throughput PICh experiment (PhD).
  \item \textbf{KARMA (iGEM 2019)} — Led the development of a modular protein degradation system; awarded Gold Medal.
\end{itemize}
\vspace{2pt}
\noindent\rule{\textwidth}{0.4pt}
% =================== EDUCATION ===================
\section*{Education}
\textbf{PhD in Genetics, Epigenetics and Cellular Determinism} \hfill 2021--2025\hspace{0.3em}\\
Université de Montpellier --- Institut de Génétique Humaine (CNRS), France

\textbf{MSc in Genetics \& Cellular Biology} \hfill 2019--2021\\
Université de Montpellier, France

\textbf{BSc in Molecular Biology} \hfill 2016--2019\\
Université de Montpellier, France\\
\vspace{2pt}
\noindent\rule{\textwidth}{0.4pt}
% =================== SELECTED METHODS \& SKILLS ===================
\section*{Skills}
\begin{itemize}
    \item \textbf{Computational:} Python (NumPy, pandas, TensorFlow, PyTorch), R (tidyverse, ggplot2, Bioconductor), Bash, Image analysis (CellProfiler, Fiji), Computer vision, Machine \& Deep Learning, Statistics, \ldots.
    \item \textbf{Experimental:} CRISPR/Cas9, ChIP, PICh, FISH (nuclei/metaphase), 3D-SIM microscopy, Proteomics, RT-qPCR, Immunofluorescence, \ldots.
    \item \textbf{Data Science:} Multi-omics integration, statistical modeling (ANOVA, linear models, clustering/PCA), reproducible pipelines, documentation, \ldots.
    \item \textbf{Communication:} Teaching (genetics, epigenetics), oral presentations, mentoring, scientific writing.
\end{itemize}
\vspace{2pt}
\noindent\rule{\textwidth}{0.4pt}
% =================== TEACHING ===================
\section*{Teaching}
\renewcommand{\labelitemi}{\faChalkboardTeacher}
\begin{itemize}
    \item Undergraduate teaching (Genetics \& Epigenetics): lectures and practicals for 2nd–3rd year students.
    \item Workshops and seminars on computational biology: taught introductory sessions on programming (Python, R), statistics, machine learning, and artificial intelligence for life scientists.
\end{itemize}
\renewcommand{\labelitemi}{\oldlabelitemi}
%\vspace{2pt}
%\noindent\rule{\textwidth}{0.4pt}
% =================== Leadership \& Community Involvement ===================
%\section*{Leadership \& Community Involvement}
%\textbf{Co-founder and Coordinator – Young Scientist Club (YSC), IGH Montpellier} \hfill 2023–Present\\
%Co-organized the institute’s Young Scientist Club; created a monthly forum fostering collaboration and knowledge exchange among PhD students, postdocs, and engineers; strengthening the institute’s research community.\\
%\vspace{2pt}
%\noindent\rule{\textwidth}{0.4pt}
%\vspace{1pt}
% =================== AWARDS ===================
%\section*{Awards \& Honors}
%iGEM Montpellier (2019): \textbf{Gold medal}.\\
%\vspace{2pt}
%\noindent\rule{\textwidth}{0.4pt}
%\vspace{2pt}
% =================== LANGUAGES ===================
%\section*{Languages}
%French (native), English (C2), Japanese (B2).\\
%\vspace{2pt}
%\noindent\rule{\textwidth}{0.4pt}
% =================== REFERENCES ===================
%\section*{References}
%Available upon request.

\end{document}
